
\documentclass{beamer}
%\documentclass[handout]{beamer}
\usepackage{pgfpages}
%\setbeameroption{show notes on second screen=left}
%\setbeameroption{hide notes}
%\setbeameroption{show only  notes}

\usepackage{xcolor}
\usepackage[utf8]{inputenc}
\usepackage[francais]{babel}
\usepackage{multimedia}
%comment
\setbeamertemplate{navigation symbols}{} 
 \usetheme{Warsaw}
\setbeamercovered{highly dynamic}

\usepackage{multimedia}

%\usepackage{pgfpages}
%\mode<handout>{\pgfpagesuselayout{4 on 1}[a4paper,landscape,border shrink=5mm]}
%\mode<handout>{\pgfpagesuselayout{8 on 1}[a4paper,border shrink=5mm]}
%\mode<handout>{\setbeamercolor{background canvas}{bg=black!5}}


\pgfdeclareimage[height=0.9\textheight]{ask}{images/ask}
\pgfdeclareimage[height=0.9\textheight]{github}{images/github}
\pgfdeclareimage[height=0.9\textheight]{qrcode}{images/qrcode}
\pgfdeclareimage[height=0.9\textheight]{githubconso1}{images/githubconso1}
\pgfdeclareimage[height=0.9\textheight]{fichierssvg}{images/fichierssvg}
\pgfdeclareimage[height=0.9\textheight]{githubpro1}{images/githubpro1}
\pgfdeclareimage[height=0.85\textheight]{astrolabe1}{images/astrolabe1}
\pgfdeclareimage[height=0.85\textheight]{astrolabe2}{images/astrolabe2}
\pgfdeclareimage[height=0.9\textheight]{davant}{images/descriptionavant.png}
\pgfdeclareimage[height=0.9\textheight]{tympanannote}{images/tympan44annote}
\pgfdeclareimage[height=0.9\textheight]{reteannote}{images/reteannote}
\pgfdeclareimage[height=0.9\textheight]{10octobre1}{images/10octobre1}
\pgfdeclareimage[height=0.9\textheight]{10octobre2}{images/10octobre2}
\pgfdeclareimage[height=0.9\textheight]{10octobre3}{images/10octobre3}
\pgfdeclareimage[height=0.9\textheight]{10octobre4}{images/10octobre4}
\pgfdeclareimage[height=0.9\textheight]{10octobre5}{images/10octobre5}
\pgfdeclareimage[height=0.9\textheight]{10octobre6}{images/10octobre6}
\pgfdeclareimage[height=0.9\textheight]{facearriere}{images/facearriere}
\pgfdeclareimage[height=0.9\textheight]{arriereannote}{images/arriereannote}
\pgfdeclareimage[width=1\textwidth]{equation}{images/equation}
\pgfdeclareimage[height=0.9\textheight]{analemme}{images/analemme}
\pgfdeclareimage[height=0.9\textheight]{exempleequation}{images/exempleequation}
\pgfdeclareimage[width=0.9\textwidth]{alidade}{images/alidade}
\pgfdeclareimage[height=0.9\textheight]{arriere}{images/arriere}
\pgfdeclareimage[width=0.9\textwidth]{ostenseur}{images/ostenseur}
\pgfdeclareimage[height=0.8\textheight]{tympan50}{images/tympan50}
\pgfdeclareimage[height=0.9\textheight]{retedebase}{images/retedebase}
\pgfdeclareimage[height=0.9\textheight]{retefinal}{images/retefinal}
\pgfdeclareimage[height=0.75\textheight]{tympan39probleme}{images/tympan39probleme}
\pgfdeclareimage[height=0.9\textheight]{tympan30probleme}{images/tympan30probleme}
\pgfdeclareimage[height=0.9\textheight]{chou}{images/chou}
\pgfdeclareimage[height=0.9\textheight]{fougere}{images/fougere}
\pgfdeclareimage[height=0.9\textheight]{fern}{images/fern}
\pgfdeclareimage[height=0.9\textheight]{foug}{images/foug}
\pgfdeclareimage[width=\textwidth]{cantor}{images/cantor}
\pgfdeclareimage[height=0.9\textheight]{koch}{images/koch}
\pgfdeclareimage[height=0.9\textheight]{hilbert}{images/hilbert}
\pgfdeclareimage[height=0.8\textheight]{mandelbrot}{images/mandelbrot}
\pgfdeclareimage[height=0.8\textheight]{brooksmatelski1978}{images/Mandel}
\pgfdeclareimage[height=0.9\textheight]{m1}{images/m1}
\pgfdeclareimage[height=0.9\textheight]{m2}{images/m2}
\pgfdeclareimage[height=0.9\textheight]{m3}{images/m3}
\pgfdeclareimage[height=0.9\textheight]{m4}{images/m4}
\pgfdeclareimage[height=0.9\textheight]{bulb1}{images/bulb1}
\pgfdeclareimage[height=0.9\textheight]{bulb2}{images/bulb2}
\pgfdeclareimage[height=0.9\textheight]{bulb3}{images/bulb3}
\pgfdeclareimage[height=0.8\textheight]{denevan}{images/denevan}



\newtheorem{fdefinition}[theorem]{Définition}
\newtheorem{Remarque}[theorem]{Remarque}
\newtheorem{Remarques}[theorem]{Remarques}
%\newtheorem{Theoreme}[theorem]{Théorème}
\newtheorem{notation}[theorem]{Notation}
  \title[Construire un astrolabe]{Construire un astrolabe}
  
  \author{Francesco De Comité}\institute{Ex-Université de Lille \\ Faculté des Sciences et Technologies}
\titlegraphic{\includegraphics[width=3.5cm]{images/Logo_fabricarium.png}\hfill}

%\titlegraphic{\ \vfill\vfill\hfill\includegraphics[width=3.5cm]{images/Logo_fabricarium.png}\hfill\hfill}


%\logo{\includegraphics[height=0.5cm]{lille1.png}}

\setbeamertemplate{part page}{
        \begin{beamercolorbox}[sep=8pt,center,wd=\textwidth]{part title}
            \usebeamerfont{part title}\insertpart\par
        \end{beamercolorbox}
}


%%%%%%%%%%%%%%%

\setbeamertemplate{part page}{
        \begin{beamercolorbox}[sep=8pt,center,wd=\textwidth, rounded=true,shadow=true]{part title}
            \usebeamerfont{part title}\insertpart\par
        \end{beamercolorbox}
}

%\definecolor{Universite}{RGB}{181,0,103}
\definecolor{Espace}{RGB}{0,37,115}

   \setbeamercolor{item projected}{bg=Espace}

\setbeamercolor{title}{bg=Espace}
\setbeamercolor{part title}{bg=Espace}
\setbeamercolor{frametitle}{bg=Espace}
%\setbeamercolor{structure}{bg=Universite}
\setbeamercolor{block title}{bg=Espace}
%\setbeamercolor{title}{fg=white,bg=!181!0!103}

\setbeamercolor{section in head/foot}{bg=Espace}

\makeatletter
\setbeamertemplate{footline}
{
  \leavevmode%
  \hbox{%
  \begin{beamercolorbox}[wd=.25\paperwidth,ht=2.25ex,dp=1ex,center]{section in head/foot}%
    \usebeamerfont{author in head/foot}\insertshortauthor
  \end{beamercolorbox}%
  \begin{beamercolorbox}[wd=.60\paperwidth,ht=2.25ex,dp=1ex,center]{section in head/foot}%
    \usebeamerfont{title in head/foot}\insertshorttitle
  \end{beamercolorbox}%
  \begin{beamercolorbox}[wd=.15\paperwidth,ht=2.25ex,dp=1ex,center]{section in head/foot}%
    %\usebeamerfont{date in head/foot}\insertshortdate{}\hspace*{2em}
     \insertframenumber/\inserttotalframenumber\hspace*{0.01ex} 
	\hfill\raisebox{-0.05cm}{\includegraphics[width=1.1cm]{images/Logo_fabricarium.png}}
  \end{beamercolorbox}}%
  \vskip0pt%
}
\makeatother


%%%%%%%%%%%%%%

\begin{document}

%%%%%%%%%%%%%%%%%%%%%%%%%%%%
 \begin{frame}
\titlepage
\note[item]{bla bla}
   \end{frame}
%%%%%%%%%%%%%%%%%%%%%%%%%%%%%%%

%%%%%%%%%%%%%%%%%%%%%%%%%%%%%%%%%%%
  
 \begin{frame}\frametitle{Généralités}

\begin{block}{Qui suis-je ? }
\begin{itemize}
\item Maître de conférences en informatique (retraité) \\ à l'Université de Lille
\item Enseignements en programmation, maths pour l'informatique, logique, intelligence artificielle. 
\item Recherche en IA, puis en Art, Mathématiques, Informatique. 
\item \begin{block}{Donner une réalité à des objets mathématiques. }
	\begin{itemize}
	\item Logiciels adéquats (Blender, Rhino, Inkscape)
	\item Langages de description (SVF, STL)
	\item Outils de conception : imprimantes 3D, découpe laser, atelier {\it ``standard''}.
	\end{itemize}
	\end{block}
\end{itemize}
\end{block}

   \end{frame}


%%%%%%%%%%%%%%%%%%%%%%%%%%%%%%%%%%%
  
 \begin{frame}\frametitle{Pourquoi construire un astrolabe ?}
 
 \begin{block}{Parce qu'on peut !}
 \begin{itemize}
 \item Des documents très bien écrits (Michel Dumas).
 \item Des logiciels.
 \item Des machines et des lieux : fablabs. 
 \item Des gens pour aider. 
 \end{itemize}
 
 
 
 \end{block}
 
\end{frame}


\begin{frame}\frametitle{N'hésitez pas à  m'interrompre !}
\begin{center}
\pgfuseimage{ask}
\end{center}
{\hfill \small Créé grace à Dall-E}
   \end{frame}

   \begin{frame}\frametitle{With a little help from my friends}
\begin{block}{Des documents très utiles}
\begin{itemize}
\item {\em Réalisation d'un astrolabe} Michel Dumas. 
\item {\em Les secrets de l'astrolabe} Ya{\"e}l Nazé.
\item Ces documents sont dans le dépôt GIT. 
\end{itemize}
\end{block}

   \end{frame}
      
   

\begin{frame}\frametitle{Toutes les infos dans le GIT}
\begin{block}{GIT minimal}
\begin{itemize}
\item GIT (Global Information tracker) permet un travail coopératif et le suivi de
versions de projets (informatiques). 
\item GIT permet de rendre des fichiers, programmes, etc... disponibles pour plusieurs utilisateurs.
\item Utilisé aussi par les {\it makers} pour diffuser leur projet. 
\item Permet un retour des utilisateurs (critiques, améliorations \dots)
\end{itemize}
\end{block}
   \end{frame}
   


\begin{frame}\frametitle{GIT en pratique -1}
{\small \tt https://github.com/francescodecomite/Construire-un-astrolabe}
{\tt https://tinyurl.com/39dbhsa5}
\begin{center}
\pgfuseimage{qrcode}
\end{center}
\vfill
\ 

\end{frame}

\begin{frame}\frametitle{GIT en pratique -- Simples consommateurs }
Simples consommateurs
\begin{center}
\hglue -1cm \pgfuseimage{githubconso1}
\end{center}

\end{frame}

\begin{frame}\frametitle{GIT en pratique -- Simples consommateurs}
\begin{center}
\hglue -1cm \pgfuseimage{fichierssvg}
\end{center}

   \end{frame}
   
   \begin{frame}\frametitle{GIT en pratique -- Pour les pros}
\begin{center}
\hglue -2cm \pgfuseimage{githubpro1}
\end{center}

   \end{frame}
   
   
   


\begin{frame}\frametitle{Le déroulé}
\begin{block}{Le plan}
\begin{itemize}
\item C'est quoi un astrolabe ?
\item Les différents composants. 
\item Qu'est-ce qu'on peut faire avec ? 
\item Comment on calcule les différentes parties ? 
\item Comment on programme ? 
\item Comment on construit ?

\end{itemize}
\end{block}

   \end{frame}
   
\begin{frame}\frametitle{Un exemple d'astrolabe : la face avant}
\begin{center}
\pgfuseimage{astrolabe1}
\end{center}
{\hfill Crédits : Landesmuseum Württemberg, Stuttgart}



   \end{frame}
   
   \begin{frame}\frametitle{Un exemple d'astrolabe : la face arrière}
\begin{center}
\pgfuseimage{astrolabe2}
\end{center}
{\hfill Crédits : Landesmuseum Württemberg, Stuttgart}
   \end{frame}
   

\begin{frame}\frametitle{La face avant en détail}
\begin{center}
\pgfuseimage{davant}
\end{center}

   \end{frame}
 
   
\begin{frame}\frametitle{Face avant. Le tympan}
\begin{center}
\pgfuseimage{tympanannote}
\end{center}
   \end{frame}


\begin{frame}\frametitle{Face avant. Le Rete}
\begin{center}
\pgfuseimage{reteannote}
\end{center}
   \end{frame} 



\begin{frame}\frametitle{La face arrière en détail}
\begin{center}
\pgfuseimage{arriereannote}
\end{center}
   \end{frame}
   
   \begin{frame}\frametitle{Face arrière complète}
\begin{center}
\pgfuseimage{facearriere}
\end{center}
   \end{frame}
   
\begin{frame}\frametitle{Un survol des fonctions}
\begin{block}{Face avant}
\begin{itemize}
\item Carte du ciel.
\item Quels sont l'heure et l'azimuth du lever de Soleil à une date donnée ? 
\item Quelle est la hauteur maximale du Soleil à une latitude donnée ? 
\item A quelle heure un étoile atteint-elle une altitude donnée, connaissant le jour et la latitude ? 
\item Que valent l'ascension droite et la déclinaison d'un astre ? 
\item {\em Pour ces questions, cf Y. Nazé}
\end{itemize}


\end{block}

   \end{frame}
   \begin{frame}\frametitle{Carte du ciel}
\begin{center}
\pgfuseimage{10octobre1}
\end{center}
   \end{frame}
   
   
      \begin{frame}\frametitle{Carte du ciel}
\begin{center}
\pgfuseimage{10octobre2}
\end{center}
   \end{frame}
 
 
    
        \begin{frame}\frametitle{Carte du ciel}
\begin{center}
\pgfuseimage{10octobre6}
\end{center}
   \end{frame}
   
   
\begin{frame}\frametitle{Carte du ciel : vérification}
\begin{center}
\pgfuseimage{10octobre5}
\end{center}
   \end{frame}
   \begin{frame}\frametitle{Un survol des fonctions}
   \begin{block}{La face arrière}
   \begin{itemize}
   \item Hauteur d'un bâtiment.
   \item Calcul des heures inégales. 
   \item Equation du temps
   \item Hauteur d'un astre au-dessus de l'horizon. 
   
   \end{itemize}
   
   \end{block}

   \end{frame}
   
\begin{frame}\frametitle{Ma préférée}
\begin{block}{L'équation du temps}
\begin{itemize}
\item A midi, il est rarement midi\dots
\begin{itemize}
\item Parce que l'orbite de la terre est une ellipse. 
\item Parce la terre est inclinée sur son orbite.
\end{itemize}
\item Le soleil peut être en avance ou en retard de 15 minutes sur son horaire. 
\end{itemize}

\end{block}

   \end{frame}
   \begin{frame}\frametitle{Equation du temps}
\begin{center}
\pgfuseimage{equation}
\end{center}
\hfill \small Crédit Wikipédia
   \end{frame}
   
\begin{frame}\frametitle{Equation du temps}
\begin{center}
\pgfuseimage{analemme}
\end{center}

\hfill \small Crédit Rhodri Evans
   \end{frame}
   \begin{frame}\frametitle{Equation du temps}
\begin{center}
\pgfuseimage{exempleequation}
\end{center}
   \end{frame}
   
   
   \begin{frame}[fragile]\frametitle{Programmer l'astrolabe : le format SVG}
\begin{block}{Primitives}
\begin{itemize}
\item \textcolor{red}{Un cercle :} 

\begin{verbatim}
<circle cx="250.0" cy="250.0"  r="7142.85" 
fill="none" stroke="black"  
stroke-width="0.5" />
\end{verbatim}
\item\textcolor{red}{Une ligne :}
\begin{verbatim}
<line x1="250" y1="250"  x2="121.7" y2="501.9"  
stroke="black"  stroke-width="0.5" />

\end{verbatim}
\item\textcolor{red}{Un arc de cercle :}
\begin{verbatim}
<path d="
M250.0 250.0
A1000 1000 0 0 0 1200 3000"  id=" " fill="none"  
stroke="black"   
stroke-width="0.5" />
\end{verbatim}

\end{itemize}

\end{block}
   \end{frame}

   



\begin{frame}[fragile]\frametitle{Construire l'astrolabe}

\begin{verbatim}
    

if __name__=="__main__":
   #alidade(R/2)  
   #ostenseur(R/2)
   #dos(R/2)
   tympan(R/2)
   tympanseul(R/2)
   #rete(R/2)
  
  

\end{verbatim}

   \end{frame}
   
   \begin{frame}\frametitle{Les fichiers prêts à découper.  Alidade}
\begin{center}
\pgfuseimage{alidade}
\end{center}
   \end{frame}
   
\begin{frame}\frametitle{Les fichiers prêts à découper.  Ostenseur}
\begin{center}
\pgfuseimage{ostenseur}
\end{center}
   \end{frame}
   
\begin{frame}\frametitle{Les fichiers prêts à découper.  Tympan50}
\begin{center}
\pgfuseimage{tympan50}
\end{center}
   \end{frame}
   

\begin{frame}\frametitle{Les fichiers prêts à découper.  Rete de base}
\begin{center}
\pgfuseimage{retedebase}
\end{center}
   \end{frame}
   
\begin{frame}\frametitle{Les fichiers prêts à découper.  Rete final}
\begin{center}
\pgfuseimage{retefinal}
\end{center}
   \end{frame}
   
\begin{frame}\frametitle{Le programme}
\begin{block}{La chaîne de fabrication}
\begin{itemize}
\item La découpeuse laser reçoit des fichiers SVG (Scalable Vector Graphics)
\item Ces fichiers lui sont envoyés depuis un logiciel de dessin (Inkscape, Illustrator,\dots)
\item Les fichiers SVG sont des fichiers texte qui contiennent les commandes de dessin
(dessine une ligne, un cercle \dots)
\item Le programme python {\tt astrolabe.py} construit un fichier SVG pour chaque partie de l'astrolabe. 
\end{itemize}

\end{block}
   \end{frame}
   

\begin{frame}\frametitle{Limites}
\begin{itemize}
\item Le programme ne dessine pas les {\bf Trônes}
\item Le {\bf rete} final peut être construit à la main.
\item Pas grave, les fichiers SVG sont fournis. 
\item Il faut être un peu attentif lorsqu'on crée de nouveaux tympans (voir détails ci-dessous). 
\end{itemize}

   \end{frame}
   
\begin{frame}\frametitle{Une nouvelle latitude}
\begin{center}
\pgfuseimage{tympan39probleme}
\end{center}
Les modifications à apporter sont toutes au même endroit dans le programme (et c'est pas dur !).
   \end{frame}
   
\begin{frame}\frametitle{Une autre nouvelle latitude}
\begin{center}
\pgfuseimage{tympan30probleme}
\end{center}
   \end{frame}
   

\begin{frame}\frametitle{Conclusions}
\begin{block}{Programme}
\begin{itemize}

\item Le programme marche souvent ;-)
\item Certains problèmes sont faciles à régler. 
\item D'autres sont plus ardus, mais maintenant je les partage avec tout le monde, dont des programmeurs plus forts que moi. 
\item GIT permet de communiquer et coopérer sur le programme. 

\end{itemize}
\end{block}

\begin{block}{Rendre le bricolage scientifique accessible}
\begin{itemize}
\item Diffusion du code source et des algorithmes (merci Michel Dumas !)
\item Utilisation de logiciels gratuits (Python, Inkscape)
\item Mise à disposition de machines (merci le Fabricarium !)
\end{itemize}
\end{block}

   \end{frame}
   
   
   \end{document}
   
\begin{frame}\frametitle{}

   \end{frame}
   
\begin{frame}\frametitle{}

   \end{frame}
   


   
   \end{document}
\begin{frame}\frametitle{}

   \end{frame}
   
\begin{frame}\frametitle{}

   \end{frame}
   
\begin{frame}\frametitle{}

   \end{frame}
   

